\documentclass{article} % For LaTeX2e
\usepackage{iclr2026_conference,times}

% Optional math commands from https://github.com/goodfeli/dlbook_notation.
%%%%% NEW MATH DEFINITIONS %%%%%

\usepackage{amsmath,amsfonts,bm}

% Mark sections of captions for referring to divisions of figures
\newcommand{\figleft}{{\em (Left)}}
\newcommand{\figcenter}{{\em (Center)}}
\newcommand{\figright}{{\em (Right)}}
\newcommand{\figtop}{{\em (Top)}}
\newcommand{\figbottom}{{\em (Bottom)}}
\newcommand{\captiona}{{\em (a)}}
\newcommand{\captionb}{{\em (b)}}
\newcommand{\captionc}{{\em (c)}}
\newcommand{\captiond}{{\em (d)}}

% Highlight a newly defined term
\newcommand{\newterm}[1]{{\bf #1}}


% Figure reference, lower-case.
\def\figref#1{figure~\ref{#1}}
% Figure reference, capital. For start of sentence
\def\Figref#1{Figure~\ref{#1}}
\def\twofigref#1#2{figures \ref{#1} and \ref{#2}}
\def\quadfigref#1#2#3#4{figures \ref{#1}, \ref{#2}, \ref{#3} and \ref{#4}}
% Section reference, lower-case.
\def\secref#1{section~\ref{#1}}
% Section reference, capital.
\def\Secref#1{Section~\ref{#1}}
% Reference to two sections.
\def\twosecrefs#1#2{sections \ref{#1} and \ref{#2}}
% Reference to three sections.
\def\secrefs#1#2#3{sections \ref{#1}, \ref{#2} and \ref{#3}}
% Reference to an equation, lower-case.
\def\eqref#1{equation~\ref{#1}}
% Reference to an equation, upper case
\def\Eqref#1{Equation~\ref{#1}}
% A raw reference to an equation---avoid using if possible
\def\plaineqref#1{\ref{#1}}
% Reference to a chapter, lower-case.
\def\chapref#1{chapter~\ref{#1}}
% Reference to an equation, upper case.
\def\Chapref#1{Chapter~\ref{#1}}
% Reference to a range of chapters
\def\rangechapref#1#2{chapters\ref{#1}--\ref{#2}}
% Reference to an algorithm, lower-case.
\def\algref#1{algorithm~\ref{#1}}
% Reference to an algorithm, upper case.
\def\Algref#1{Algorithm~\ref{#1}}
\def\twoalgref#1#2{algorithms \ref{#1} and \ref{#2}}
\def\Twoalgref#1#2{Algorithms \ref{#1} and \ref{#2}}
% Reference to a part, lower case
\def\partref#1{part~\ref{#1}}
% Reference to a part, upper case
\def\Partref#1{Part~\ref{#1}}
\def\twopartref#1#2{parts \ref{#1} and \ref{#2}}

\def\ceil#1{\lceil #1 \rceil}
\def\floor#1{\lfloor #1 \rfloor}
\def\1{\bm{1}}
\newcommand{\train}{\mathcal{D}}
\newcommand{\valid}{\mathcal{D_{\mathrm{valid}}}}
\newcommand{\test}{\mathcal{D_{\mathrm{test}}}}

\def\eps{{\epsilon}}


% Random variables
\def\reta{{\textnormal{$\eta$}}}
\def\ra{{\textnormal{a}}}
\def\rb{{\textnormal{b}}}
\def\rc{{\textnormal{c}}}
\def\rd{{\textnormal{d}}}
\def\re{{\textnormal{e}}}
\def\rf{{\textnormal{f}}}
\def\rg{{\textnormal{g}}}
\def\rh{{\textnormal{h}}}
\def\ri{{\textnormal{i}}}
\def\rj{{\textnormal{j}}}
\def\rk{{\textnormal{k}}}
\def\rl{{\textnormal{l}}}
% rm is already a command, just don't name any random variables m
\def\rn{{\textnormal{n}}}
\def\ro{{\textnormal{o}}}
\def\rp{{\textnormal{p}}}
\def\rq{{\textnormal{q}}}
\def\rr{{\textnormal{r}}}
\def\rs{{\textnormal{s}}}
\def\rt{{\textnormal{t}}}
\def\ru{{\textnormal{u}}}
\def\rv{{\textnormal{v}}}
\def\rw{{\textnormal{w}}}
\def\rx{{\textnormal{x}}}
\def\ry{{\textnormal{y}}}
\def\rz{{\textnormal{z}}}

% Random vectors
\def\rvepsilon{{\mathbf{\epsilon}}}
\def\rvtheta{{\mathbf{\theta}}}
\def\rva{{\mathbf{a}}}
\def\rvb{{\mathbf{b}}}
\def\rvc{{\mathbf{c}}}
\def\rvd{{\mathbf{d}}}
\def\rve{{\mathbf{e}}}
\def\rvf{{\mathbf{f}}}
\def\rvg{{\mathbf{g}}}
\def\rvh{{\mathbf{h}}}
\def\rvu{{\mathbf{i}}}
\def\rvj{{\mathbf{j}}}
\def\rvk{{\mathbf{k}}}
\def\rvl{{\mathbf{l}}}
\def\rvm{{\mathbf{m}}}
\def\rvn{{\mathbf{n}}}
\def\rvo{{\mathbf{o}}}
\def\rvp{{\mathbf{p}}}
\def\rvq{{\mathbf{q}}}
\def\rvr{{\mathbf{r}}}
\def\rvs{{\mathbf{s}}}
\def\rvt{{\mathbf{t}}}
\def\rvu{{\mathbf{u}}}
\def\rvv{{\mathbf{v}}}
\def\rvw{{\mathbf{w}}}
\def\rvx{{\mathbf{x}}}
\def\rvy{{\mathbf{y}}}
\def\rvz{{\mathbf{z}}}

% Elements of random vectors
\def\erva{{\textnormal{a}}}
\def\ervb{{\textnormal{b}}}
\def\ervc{{\textnormal{c}}}
\def\ervd{{\textnormal{d}}}
\def\erve{{\textnormal{e}}}
\def\ervf{{\textnormal{f}}}
\def\ervg{{\textnormal{g}}}
\def\ervh{{\textnormal{h}}}
\def\ervi{{\textnormal{i}}}
\def\ervj{{\textnormal{j}}}
\def\ervk{{\textnormal{k}}}
\def\ervl{{\textnormal{l}}}
\def\ervm{{\textnormal{m}}}
\def\ervn{{\textnormal{n}}}
\def\ervo{{\textnormal{o}}}
\def\ervp{{\textnormal{p}}}
\def\ervq{{\textnormal{q}}}
\def\ervr{{\textnormal{r}}}
\def\ervs{{\textnormal{s}}}
\def\ervt{{\textnormal{t}}}
\def\ervu{{\textnormal{u}}}
\def\ervv{{\textnormal{v}}}
\def\ervw{{\textnormal{w}}}
\def\ervx{{\textnormal{x}}}
\def\ervy{{\textnormal{y}}}
\def\ervz{{\textnormal{z}}}

% Random matrices
\def\rmA{{\mathbf{A}}}
\def\rmB{{\mathbf{B}}}
\def\rmC{{\mathbf{C}}}
\def\rmD{{\mathbf{D}}}
\def\rmE{{\mathbf{E}}}
\def\rmF{{\mathbf{F}}}
\def\rmG{{\mathbf{G}}}
\def\rmH{{\mathbf{H}}}
\def\rmI{{\mathbf{I}}}
\def\rmJ{{\mathbf{J}}}
\def\rmK{{\mathbf{K}}}
\def\rmL{{\mathbf{L}}}
\def\rmM{{\mathbf{M}}}
\def\rmN{{\mathbf{N}}}
\def\rmO{{\mathbf{O}}}
\def\rmP{{\mathbf{P}}}
\def\rmQ{{\mathbf{Q}}}
\def\rmR{{\mathbf{R}}}
\def\rmS{{\mathbf{S}}}
\def\rmT{{\mathbf{T}}}
\def\rmU{{\mathbf{U}}}
\def\rmV{{\mathbf{V}}}
\def\rmW{{\mathbf{W}}}
\def\rmX{{\mathbf{X}}}
\def\rmY{{\mathbf{Y}}}
\def\rmZ{{\mathbf{Z}}}

% Elements of random matrices
\def\ermA{{\textnormal{A}}}
\def\ermB{{\textnormal{B}}}
\def\ermC{{\textnormal{C}}}
\def\ermD{{\textnormal{D}}}
\def\ermE{{\textnormal{E}}}
\def\ermF{{\textnormal{F}}}
\def\ermG{{\textnormal{G}}}
\def\ermH{{\textnormal{H}}}
\def\ermI{{\textnormal{I}}}
\def\ermJ{{\textnormal{J}}}
\def\ermK{{\textnormal{K}}}
\def\ermL{{\textnormal{L}}}
\def\ermM{{\textnormal{M}}}
\def\ermN{{\textnormal{N}}}
\def\ermO{{\textnormal{O}}}
\def\ermP{{\textnormal{P}}}
\def\ermQ{{\textnormal{Q}}}
\def\ermR{{\textnormal{R}}}
\def\ermS{{\textnormal{S}}}
\def\ermT{{\textnormal{T}}}
\def\ermU{{\textnormal{U}}}
\def\ermV{{\textnormal{V}}}
\def\ermW{{\textnormal{W}}}
\def\ermX{{\textnormal{X}}}
\def\ermY{{\textnormal{Y}}}
\def\ermZ{{\textnormal{Z}}}

% Vectors
\def\vzero{{\bm{0}}}
\def\vone{{\bm{1}}}
\def\vmu{{\bm{\mu}}}
\def\vtheta{{\bm{\theta}}}
\def\va{{\bm{a}}}
\def\vb{{\bm{b}}}
\def\vc{{\bm{c}}}
\def\vd{{\bm{d}}}
\def\ve{{\bm{e}}}
\def\vf{{\bm{f}}}
\def\vg{{\bm{g}}}
\def\vh{{\bm{h}}}
\def\vi{{\bm{i}}}
\def\vj{{\bm{j}}}
\def\vk{{\bm{k}}}
\def\vl{{\bm{l}}}
\def\vm{{\bm{m}}}
\def\vn{{\bm{n}}}
\def\vo{{\bm{o}}}
\def\vp{{\bm{p}}}
\def\vq{{\bm{q}}}
\def\vr{{\bm{r}}}
\def\vs{{\bm{s}}}
\def\vt{{\bm{t}}}
\def\vu{{\bm{u}}}
\def\vv{{\bm{v}}}
\def\vw{{\bm{w}}}
\def\vx{{\bm{x}}}
\def\vy{{\bm{y}}}
\def\vz{{\bm{z}}}

% Elements of vectors
\def\evalpha{{\alpha}}
\def\evbeta{{\beta}}
\def\evepsilon{{\epsilon}}
\def\evlambda{{\lambda}}
\def\evomega{{\omega}}
\def\evmu{{\mu}}
\def\evpsi{{\psi}}
\def\evsigma{{\sigma}}
\def\evtheta{{\theta}}
\def\eva{{a}}
\def\evb{{b}}
\def\evc{{c}}
\def\evd{{d}}
\def\eve{{e}}
\def\evf{{f}}
\def\evg{{g}}
\def\evh{{h}}
\def\evi{{i}}
\def\evj{{j}}
\def\evk{{k}}
\def\evl{{l}}
\def\evm{{m}}
\def\evn{{n}}
\def\evo{{o}}
\def\evp{{p}}
\def\evq{{q}}
\def\evr{{r}}
\def\evs{{s}}
\def\evt{{t}}
\def\evu{{u}}
\def\evv{{v}}
\def\evw{{w}}
\def\evx{{x}}
\def\evy{{y}}
\def\evz{{z}}

% Matrix
\def\mA{{\bm{A}}}
\def\mB{{\bm{B}}}
\def\mC{{\bm{C}}}
\def\mD{{\bm{D}}}
\def\mE{{\bm{E}}}
\def\mF{{\bm{F}}}
\def\mG{{\bm{G}}}
\def\mH{{\bm{H}}}
\def\mI{{\bm{I}}}
\def\mJ{{\bm{J}}}
\def\mK{{\bm{K}}}
\def\mL{{\bm{L}}}
\def\mM{{\bm{M}}}
\def\mN{{\bm{N}}}
\def\mO{{\bm{O}}}
\def\mP{{\bm{P}}}
\def\mQ{{\bm{Q}}}
\def\mR{{\bm{R}}}
\def\mS{{\bm{S}}}
\def\mT{{\bm{T}}}
\def\mU{{\bm{U}}}
\def\mV{{\bm{V}}}
\def\mW{{\bm{W}}}
\def\mX{{\bm{X}}}
\def\mY{{\bm{Y}}}
\def\mZ{{\bm{Z}}}
\def\mBeta{{\bm{\beta}}}
\def\mPhi{{\bm{\Phi}}}
\def\mLambda{{\bm{\Lambda}}}
\def\mSigma{{\bm{\Sigma}}}

% Tensor
\DeclareMathAlphabet{\mathsfit}{\encodingdefault}{\sfdefault}{m}{sl}
\SetMathAlphabet{\mathsfit}{bold}{\encodingdefault}{\sfdefault}{bx}{n}
\newcommand{\tens}[1]{\bm{\mathsfit{#1}}}
\def\tA{{\tens{A}}}
\def\tB{{\tens{B}}}
\def\tC{{\tens{C}}}
\def\tD{{\tens{D}}}
\def\tE{{\tens{E}}}
\def\tF{{\tens{F}}}
\def\tG{{\tens{G}}}
\def\tH{{\tens{H}}}
\def\tI{{\tens{I}}}
\def\tJ{{\tens{J}}}
\def\tK{{\tens{K}}}
\def\tL{{\tens{L}}}
\def\tM{{\tens{M}}}
\def\tN{{\tens{N}}}
\def\tO{{\tens{O}}}
\def\tP{{\tens{P}}}
\def\tQ{{\tens{Q}}}
\def\tR{{\tens{R}}}
\def\tS{{\tens{S}}}
\def\tT{{\tens{T}}}
\def\tU{{\tens{U}}}
\def\tV{{\tens{V}}}
\def\tW{{\tens{W}}}
\def\tX{{\tens{X}}}
\def\tY{{\tens{Y}}}
\def\tZ{{\tens{Z}}}


% Graph
\def\gA{{\mathcal{A}}}
\def\gB{{\mathcal{B}}}
\def\gC{{\mathcal{C}}}
\def\gD{{\mathcal{D}}}
\def\gE{{\mathcal{E}}}
\def\gF{{\mathcal{F}}}
\def\gG{{\mathcal{G}}}
\def\gH{{\mathcal{H}}}
\def\gI{{\mathcal{I}}}
\def\gJ{{\mathcal{J}}}
\def\gK{{\mathcal{K}}}
\def\gL{{\mathcal{L}}}
\def\gM{{\mathcal{M}}}
\def\gN{{\mathcal{N}}}
\def\gO{{\mathcal{O}}}
\def\gP{{\mathcal{P}}}
\def\gQ{{\mathcal{Q}}}
\def\gR{{\mathcal{R}}}
\def\gS{{\mathcal{S}}}
\def\gT{{\mathcal{T}}}
\def\gU{{\mathcal{U}}}
\def\gV{{\mathcal{V}}}
\def\gW{{\mathcal{W}}}
\def\gX{{\mathcal{X}}}
\def\gY{{\mathcal{Y}}}
\def\gZ{{\mathcal{Z}}}

% Sets
\def\sA{{\mathbb{A}}}
\def\sB{{\mathbb{B}}}
\def\sC{{\mathbb{C}}}
\def\sD{{\mathbb{D}}}
% Don't use a set called E, because this would be the same as our symbol
% for expectation.
\def\sF{{\mathbb{F}}}
\def\sG{{\mathbb{G}}}
\def\sH{{\mathbb{H}}}
\def\sI{{\mathbb{I}}}
\def\sJ{{\mathbb{J}}}
\def\sK{{\mathbb{K}}}
\def\sL{{\mathbb{L}}}
\def\sM{{\mathbb{M}}}
\def\sN{{\mathbb{N}}}
\def\sO{{\mathbb{O}}}
\def\sP{{\mathbb{P}}}
\def\sQ{{\mathbb{Q}}}
\def\sR{{\mathbb{R}}}
\def\sS{{\mathbb{S}}}
\def\sT{{\mathbb{T}}}
\def\sU{{\mathbb{U}}}
\def\sV{{\mathbb{V}}}
\def\sW{{\mathbb{W}}}
\def\sX{{\mathbb{X}}}
\def\sY{{\mathbb{Y}}}
\def\sZ{{\mathbb{Z}}}

% Entries of a matrix
\def\emLambda{{\Lambda}}
\def\emA{{A}}
\def\emB{{B}}
\def\emC{{C}}
\def\emD{{D}}
\def\emE{{E}}
\def\emF{{F}}
\def\emG{{G}}
\def\emH{{H}}
\def\emI{{I}}
\def\emJ{{J}}
\def\emK{{K}}
\def\emL{{L}}
\def\emM{{M}}
\def\emN{{N}}
\def\emO{{O}}
\def\emP{{P}}
\def\emQ{{Q}}
\def\emR{{R}}
\def\emS{{S}}
\def\emT{{T}}
\def\emU{{U}}
\def\emV{{V}}
\def\emW{{W}}
\def\emX{{X}}
\def\emY{{Y}}
\def\emZ{{Z}}
\def\emSigma{{\Sigma}}

% entries of a tensor
% Same font as tensor, without \bm wrapper
\newcommand{\etens}[1]{\mathsfit{#1}}
\def\etLambda{{\etens{\Lambda}}}
\def\etA{{\etens{A}}}
\def\etB{{\etens{B}}}
\def\etC{{\etens{C}}}
\def\etD{{\etens{D}}}
\def\etE{{\etens{E}}}
\def\etF{{\etens{F}}}
\def\etG{{\etens{G}}}
\def\etH{{\etens{H}}}
\def\etI{{\etens{I}}}
\def\etJ{{\etens{J}}}
\def\etK{{\etens{K}}}
\def\etL{{\etens{L}}}
\def\etM{{\etens{M}}}
\def\etN{{\etens{N}}}
\def\etO{{\etens{O}}}
\def\etP{{\etens{P}}}
\def\etQ{{\etens{Q}}}
\def\etR{{\etens{R}}}
\def\etS{{\etens{S}}}
\def\etT{{\etens{T}}}
\def\etU{{\etens{U}}}
\def\etV{{\etens{V}}}
\def\etW{{\etens{W}}}
\def\etX{{\etens{X}}}
\def\etY{{\etens{Y}}}
\def\etZ{{\etens{Z}}}

% The true underlying data generating distribution
\newcommand{\pdata}{p_{\rm{data}}}
% The empirical distribution defined by the training set
\newcommand{\ptrain}{\hat{p}_{\rm{data}}}
\newcommand{\Ptrain}{\hat{P}_{\rm{data}}}
% The model distribution
\newcommand{\pmodel}{p_{\rm{model}}}
\newcommand{\Pmodel}{P_{\rm{model}}}
\newcommand{\ptildemodel}{\tilde{p}_{\rm{model}}}
% Stochastic autoencoder distributions
\newcommand{\pencode}{p_{\rm{encoder}}}
\newcommand{\pdecode}{p_{\rm{decoder}}}
\newcommand{\precons}{p_{\rm{reconstruct}}}

\newcommand{\laplace}{\mathrm{Laplace}} % Laplace distribution

\newcommand{\E}{\mathbb{E}}
\newcommand{\Ls}{\mathcal{L}}
\newcommand{\R}{\mathbb{R}}
\newcommand{\emp}{\tilde{p}}
\newcommand{\lr}{\alpha}
\newcommand{\reg}{\lambda}
\newcommand{\rect}{\mathrm{rectifier}}
\newcommand{\softmax}{\mathrm{softmax}}
\newcommand{\sigmoid}{\sigma}
\newcommand{\softplus}{\zeta}
\newcommand{\KL}{D_{\mathrm{KL}}}
\newcommand{\Var}{\mathrm{Var}}
\newcommand{\standarderror}{\mathrm{SE}}
\newcommand{\Cov}{\mathrm{Cov}}
% Wolfram Mathworld says $L^2$ is for function spaces and $\ell^2$ is for vectors
% But then they seem to use $L^2$ for vectors throughout the site, and so does
% wikipedia.
\newcommand{\normlzero}{L^0}
\newcommand{\normlone}{L^1}
\newcommand{\normltwo}{L^2}
\newcommand{\normlp}{L^p}
\newcommand{\normmax}{L^\infty}

\newcommand{\parents}{Pa} % See usage in notation.tex. Chosen to match Daphne's book.

\DeclareMathOperator*{\argmax}{arg\,max}
\DeclareMathOperator*{\argmin}{arg\,min}

\DeclareMathOperator{\sign}{sign}
\DeclareMathOperator{\Tr}{Tr}
\let\ab\allowbreak


\usepackage{hyperref}
\usepackage{url}
\usepackage{graphicx}
\usepackage{booktabs}
\usepackage{multirow}
\usepackage{xcolor}
\usepackage{subcaption}
\usepackage{wrapfig}

% Placeholder marker for sections not written yet. Remove before submission.
\newcommand{\todo}[1]{\textcolor{red}{[TODO: #1]}}

\graphicspath{{../../figures/}{../slide-images/}}

\title{Where the Accuracy Comes From: Recipe, Resolution, \\
Parameter Placement and Prediction-Time Re-ranking \\
for Adapting ViT-B/16 to 10,000 Fine-Grained Classes}

\author{Anastasia Frances Frederica, Le Thi Hong Van, Liu Wenqi \& Xu Qianqian \\
CS604 Team 08 \\
Singapore Management University
}

\iclrfinalcopy % Show author names (course report, not an anonymous submission).

% \iclrfinalcopy keeps the author names. The style then writes
% "Published as a conference paper at ICLR 2026" from inside \maketitle,
% and a later \pagestyle{fancy} would put that line back. Install the
% course header whenever the fancy page style is applied.
\makeatletter
\let\course@origfancy\ps@@fancy
\def\ps@@fancy{%
  \course@origfancy
  \def\@oddhead{\hfill CS604 Course Project Report\hfill}%
  \def\@evenhead{\hfill CS604 Course Project Report\hfill}%
}
\pagestyle{fancy}
\let\course@origtitle\maketitle
\def\maketitle{\course@origtitle\fancyhead{}\pagestyle{fancy}}
\makeatother

\begin{document}

\maketitle

% ---------------------------------------------------------------------------
\begin{abstract}
\todo{About 200 words: problem, three RQs, key numbers (full FT 84.07\%, AF + LoRA
at 4.94\% 83.08\%, zoom + kNN +2.0 pp on the holdout without retraining, +0.93 on
Phase 2), one-line lesson.}
\end{abstract}

% ---------------------------------------------------------------------------
\section{Introduction}
\label{sec:intro}

A pretrained Vision Transformer (ViT) is rarely used as it is: it is adapted to the task
at hand. We study this adaptation in a setting that is hard on two counts. The task is
fine-grained classification over 10,000 classes with only 45 training images per class,
where closely related classes share shape, colour and texture (inter-class similarity) and
images of one class vary in pose, life stage, lighting and background (intra-class
variation). The compute is also fixed: every training run must fit in 2 hours on 48\,GB of
GPU memory, starting from one ViT-B/16 checkpoint with 86M parameters.

Two practical tensions motivate the project. First, full fine-tuning gives the model the
most freedom, but every new task then needs its own 86M-parameter copy and the memory to
train it. Parameter-efficient fine-tuning (PEFT) freezes the backbone and trains small
modules instead, typically under 5\% of the parameters. What matters is how much accuracy
that costs. Second, PEFT methods are usually compared by how many parameters they train.
Our budget, however, is set in wall-clock time, and inference cost is set by latency, so
training fewer parameters only helps if it also saves time, memory or latency. A third
question follows once a model is trained: its remaining errors are costly to fix by
retraining, so we ask whether they can be reduced more cheaply.

Our first experiment reframed the project. Fine-tuning all 86M weights for the full
budget reached 61.3\% top-1, \emph{below} a linear probe trained on frozen features
(62.5\%). Giving every weight freedom to change did not help. The limit was \emph{how}
the weights were changed: the head's starting point, the learning rate of each block and
the input resolution. We therefore treat adaptation as a joint design problem, in which
these choices and the prediction procedure count as much as the adaptation method. A
second consequence is that any comparison between full fine-tuning and PEFT needs a
strong full fine-tuning baseline. Comparing against the 61.3\% run would flatter PEFT.

We ask three research questions:
\begin{itemize}
  \item[\textbf{RQ1}] \textbf{Full fine-tuning versus PEFT under a fixed budget.} With the
  same 2-hour budget and the same prediction pipeline, how much accuracy does a PEFT
  method training under 5\% of the backbone give up relative to a well-tuned full
  fine-tune, and which costs does it actually reduce: trainable parameters, training
  memory, training time or inference latency?
  \item[\textbf{RQ2}] \textbf{Spending a fixed PEFT budget.} With the trainable share held
  below 5\%, how much does accuracy depend on the choice of method, on \emph{where} the
  parameters sit in the network, and on \emph{what kind} of capacity they add?
  \item[\textbf{RQ3}] \textbf{What limits the adapted model.} What kinds of error remain
  after adaptation, and can prediction-time steps that train nothing reduce them?
\end{itemize}

Our contributions are as follows.
\begin{itemize}
  \item \textbf{A strong full fine-tuning baseline under the budget}
  (Sections~\ref{sec:method} and~\ref{sec:results}). A prototype-initialised cosine head
  and layer-wise learning-rate decay raise full fine-tuning from 66.3\% to 75.7\% at
  224\,px. Spending the budget on resolution rather than epochs raises it to 82.5\% at
  448\,px. At equal compute, 448\,px beats 224\,px even though it fits about half as many
epochs (4.1 vs 7.7).
  \item \textbf{A controlled study of how to spend a PEFT budget.} \emph{Which} parameters
  are trained matters more than \emph{how many}. BitFit reaches 73.4\% with 0.12\% of the
  backbone, while last-block attention reaches 67.9\% with 2.75\%. Cheap biases and
  LayerNorms in every block help. Once every block can adapt, placement barely matters,
  and LoRA is the efficient way to add capacity. The resulting AdaptFormer + LoRA design
  trains 4.94\% of the backbone and finishes within 1.0 percentage point (pp) of full
  fine-tuning (83.08\% vs 84.07\%). It saves training memory (14.66 vs 17.80\,GB) but no
  training time (both complete 5.5 epochs), and is slower at inference (11.27 vs 8.47\,ms
  per image).
  \item \textbf{An error analysis that motivates two training-free extensions.} About half
  of the remaining errors confuse look-alike relatives and about half involve
  appearances unlike any training image of the class. In 64\% of errors the true class is
  still in the top five. An attention-guided zoom (a second, closer view) and kNN
  re-ranking of the top five against individual training images together add
  +2.0\,pp on our holdout. With settings fixed in advance, the kNN gain held on the
  hidden Phase~2 test set (+0.93\,pp).
  \item \textbf{An analysis of training dynamics and overfitting} for both final systems
  (Section~\ref{sec:results}). \todo{one-sentence finding once the loss and
  train-vs-holdout curves are in.}
\end{itemize}

% ---------------------------------------------------------------------------
\section{Related Work}
\label{sec:related}

\paragraph{Pretrained Vision Transformers and transfer.}
\todo{ViT \citep{dosovitskiy2021vit}, AugReg \citep{steiner2022augreg}.}

\paragraph{Full fine-tuning recipes.}
\todo{LP-FT \citep{kumar2022lpft}, layer-wise LR decay \citep{bao2022beit,he2022mae},
cosine and margin heads \citep{wang2018cosface}, class-prototype classifiers
\citep{mensink2013ncm,snell2017protonet}. An alternative is to start each class weight from
the mean feature of that class, as in weight imprinting \citep{qi2018imprinting}. Resolution \citep{touvron2019fixres}.
Pros and cons.}

\paragraph{Parameter-efficient fine-tuning.}
\todo{BitFit \citep{zaken2022bitfit}, prefix tuning and VPT \citep{li2021prefix,jia2022vpt},
adapters \citep{houlsby2019adapter}, AdaptFormer \citep{chen2022adaptformer},
LoRA \citep{hu2022lora}, DoRA \citep{liu2024dora}, AdaLoRA \citep{zhang2023adalora},
GPS \citep{zhang2024gps}, surgical FT \citep{lee2023surgical}.
Pros and cons; Table~\ref{tab:related}.}

\paragraph{Prediction-time methods.}
\todo{Test-time augmentation \citep{krizhevsky2012imagenet}; attention-guided zoom for fine-grained recognition
\citep{fu2017racnn,hu2019wsdan}; kNN and memory-based classification
\citep{nakata2022knn,wu2018instance,khandelwal2020knnlm}.}

\todo{Gap: methods are usually compared at equal parameters; we compare at equal
wall-clock and treat the head, resolution and prediction as part of the method.}

% Table 1: method | what trains | mergeable | pros | cons
% \begin{table}[t] ... \label{tab:related} \end{table}

% ---------------------------------------------------------------------------
\section{Method}
\label{sec:method}

% Each subsection: What we did -> Why (evidence or hypothesis) -> What we rejected.

\subsection{Problem setup and budget}
\paragraph{Dataset.} We use the CS604 course project dataset \citep{cs604data}, an anonymised
fine-grained image classification dataset with 10,000 classes. It provides 450,000 labelled
training images (45 per class), a validation set of 50,000 images (5 per class) and a test set
of 100,000 images (10 per class). Validation and test labels are not available to us, so both
sets are scored only on the evaluation server (Section~\ref{sec:setup}). Class names are
anonymised, so no class semantics or hierarchy is available.

\todo{ViT-B/16 (augreg2) \citep{steiner2022augreg} pretrained on ImageNet-21k \citep{deng2009imagenet}, loaded via timm \citep{rw2019timm} on PyTorch \citep{paszke2019pytorch} -- 2 h $\times$ 48 GB
(144-minute runs), metrics.}

\subsection{Overview}
\todo{Figure~\ref{fig:pipeline}: three lanes (full FT / PEFT / prediction time),
our changes highlighted.}

% \begin{figure}[t]
%   \centering
%   \includegraphics[width=\linewidth]{pipeline.pdf}
%   \caption{...}
%   \label{fig:pipeline}
% \end{figure}

\subsection{Full fine-tuning recipe}
\todo{Diagnosis; prototype cosine head \citep{mensink2013ncm,snell2017protonet} (why not LP-FT); layer-wise LR decay 0.75;
schedule and regularisation (details in Appendix). Untuned settings follow the MAE ViT-B fine-tuning recipe
\citep{he2022mae}: label smoothing 0.1 \citep{szegedy2016inception}, drop path 0.1 \citep{huang2016stochdepth},
AdamW weight decay 0.05 \citep{loshchilov2019adamw}, no decay on norms and biases \citep{he2019bagoftricks};
gradient clipping at norm 1.0 as in ViT fine-tuning \citep{dosovitskiy2021vit}.}

\subsection{Spending compute on pixels, not epochs}
\todo{Small subjects (12 vs 35 patches); trade-off against epochs.}

\subsection{PEFT design}
\todo{AF w192 + LoRA + biases/norms + cosine margin at 4.94\%; why AF is kept;
latency cost of the branch that cannot be merged.}

\subsection{Prediction-time extensions}
\todo{Zoom second view from a class activation map \citep{zhou2016cam}; kNN re-ranking, score $= \log p + 20 \times$ mean of the
top-3 cosines over the top-5 list; settings fixed on half of the classes.}

\subsection{Parameter accounting}
\todo{What counts towards the 5\% budget (head excluded).}

% ---------------------------------------------------------------------------
\section{Experimental Setup}
\label{sec:setup}
\todo{Three-stage protocol (pilot, labelled holdout, CodaBench) and why; hardware,
bf16, batch sizes, seeds, tie threshold, latency protocol.}

% ---------------------------------------------------------------------------
\section{Results and Analysis}
\label{sec:results}

\subsection{Main results}
\todo{Table: final systems on Phase 1 and Phase 2: top-1, trainable \%, VRAM,
latency, epochs, GPU-hours.}

\subsection{Full fine-tuning ablations}
\todo{LP-FT vs prototype; layer decay; resolution at equal compute; waterfall.}

\subsection{PEFT comparisons}
\todo{Method scorecard, cost table, placement, capacity type, build path.}

\subsection{Training dynamics and overfitting}
\todo{Loss curves of both finals; train vs holdout curves; overfitting vs
underfitting; diminishing epochs; loss levels are not comparable across systems.}

\subsection{Error analysis}
\todo{Example types, cause shares, error rate by flag, top-k.}

\subsection{Prediction-time extensions}
\todo{Second view, kNN variants, zoom + kNN additivity, Phase 2 transfer and cost.}

% ---------------------------------------------------------------------------
\section{Discussion and Conclusion}
\label{sec:conclusion}
\todo{Answers to RQ1--RQ3; negative results; limitations; recommendations.}

% ---------------------------------------------------------------------------
\bibliography{refs}
\bibliographystyle{iclr2026_conference}

\appendix
\section{Appendix}
\todo{LR sweep; full FT recipe; head $\times$ layer decay; PEFT tuning grids;
error-flag definitions; training logs; team contributions.}

\end{document}
